% RESEARCH_DRAFT_SHA256: 1bd7258859e019ab7aa43ff440c7474598abf4491da7a13a86a480c8a57e9301
% Layout and source-reference wrappers expanded for the web renderer.
\documentclass[11pt]{article}
\usepackage[margin=25mm]{geometry}
\usepackage{fontspec,unicode-math}
\setmainfont{Latin Modern Roman}
\setmathfont{Latin Modern Math}
\usepackage{amsmath,amssymb,mathtools,xurl,hyperref}
\hypersetup{colorlinks=true,urlcolor=blue,linkcolor=blue}
\setcounter{secnumdepth}{0}
\setlength{\parindent}{0pt}
\setlength{\parskip}{5pt}
\setlength{\emergencystretch}{3em}
\newcommand{\sref}[2]{\hyperlink{source:#1:#2}{[S#2]}}
\newcommand{\cataloguescope}[1]{\paragraph{Source record.}#1}
\newcommand{\floor}[1]{\left\lfloor#1\right\rfloor}
\begin{document}
\section*{Research-attempt provenance}
This standalone research attempt retains the complete original definition body below,
including its historical source and preparation notes. Those notes describe the
earlier catalogue preparation; the new attempt and its verification are identified
separately after that body. No original ID, catalogue status, or source statement is changed.

\section{Definition of a Causal Representation}\label{statement}

\subsection{Definitions and mathematical statement}

\paragraph{Shared notation.}Standard number systems are $\mathbb N,\mathbb Z,\mathbb Q,\mathbb R,\mathbb C$. Specialized notation and index ranges follow the cited source and the definitions below.


A representation maps observations to features. A causal model describes how interventions or structural mechanisms change the distribution, rather than recording associations alone.
\paragraph{Record scope and qualifications.}
The catalogue title is retained as metadata. Named certificates, contracts, bounds or frameworks in that title are not asserted as results of the original source extract.
The supplied registry flags this extraction as unrecoverable. It is displayed as a literal source fragment; the definitions provide context rather than a replacement formal statement.

No specific structural causal model, intervention family, or required representation property is supplied.

Source-record limitation: Source flags the statement as unrecoverable: The official report presents this as a broad epistemological framing question rather than a precise mathematical problem with stated hypotheses and conclusion.
The source text below is rendered literally where necessary; damaged formulas and composite extracts are not silently repaired into a new mathematical claim.
The supplied catalogue marks this record as unrecoverable. Its extract and terminology are documented, but a complete original problem has not been reconstructed.
\paragraph{Statement recorded in the supplied source.}
What is a causal representation? The report passage is retained verbatim as the relevant context for this extracted problem. \href{https://doi.org/10.4171/owr/2025/4}{[S2]}
\paragraph{Source status.}
The supplied catalogue labels this source record open; that metadata was checked against the supplied snapshot on 2026-10-11. The label is preserved for provenance and is not an independent verification of present open status or mathematical correctness. The catalogue literature assessment is dated 2026-08-22.
\subsection{Short English statement}
The fragment asks for a definition of causal representation.
\subsection{Sources}
\begin{itemize}
\item[S1] Ulam AI and the UnsolvedMath Contributors, UnsolvedMath: A Curated Collection of Open Mathematics Problems, catalogue entry. Catalogue curation and metadata: CC BY 4.0. The source extract is retained; definitions, English summaries and qualifications are editorial additions.
\url{https://huggingface.co/datasets/ulamai/UnsolvedMath}.
\item[S2] Overparametrization, Regularization, Identifiability and Uncertainty in Machine Learning (2025), Oberwolfach Reports, DOI 10.4171/owr/2025/4, pp. 175--208
\url{https://doi.org/10.4171/owr/2025/4}.
\item[S3] Oberwolfach Report 4/2025
\url{https://content.ems.press/assets/public/full-texts/serials/owr/22/1/14299083/online/10.4171-owr-2025-4.pdf}.
\item[S4] Scholkopf, What is a causal representation
\url{https://nips.cc/virtual/2024/108171}.
\item Source-specific Oberwolfach Reports licences, including CC BY-SA 4.0 where applicable, are retained; see THIRD\_PARTY\_NOTICES.md.
\end{itemize}
\paragraph{Preparation and review.}
The original source and UnsolvedMath contributors retain attribution for the source extract. Codex (GPT-6; exact serving revision not exposed) assisted with definitions, English summaries, scope notes and formatting. No new human mathematical authorship or independently verified solution is asserted. Catalogue identity checked against the complete supplied registry and current repository. Distinct upstream records, including repeated questions and broad or damaged extracts, are retained for complete coverage of the source CS category. A repeated record is not a new independent conjecture. Complete mathematical-equivalence and current-literature review has not been performed. Preparation and final review follow contribution prompts C and D, with source gaps disclosed. Structural and rendering checks do not establish a new theorem or open status.
\label{end:statement}

\section{Literature and context check}
The primary report, printed pp. 199--202, already gives commuting-diagram definitions involving representations, transformations, intervention maps, and pushforwards. Thus the abbreviated question is not treated as an absence of any definition. The following exact finite deterministic quotient criterion investigates one restricted part of that framework.
\section{Research attempt: Finite deterministic intervention quotients}

Fix a finite nonempty set $X$, maps $A_1,\ldots,A_k:X\to X$, and a surjection $\phi:X\to Z$. There exist maps $\overline A_i:Z\to Z$ satisfying
\[
\phi\circ A_i=\overline A_i\circ\phi
\]
if and only if every fiber equivalence is forward stable:
\[
\phi(x)=\phi(y)\quad\Longrightarrow\quad
\phi(A_i x)=\phi(A_i y)\quad\text{for every }i.
\]
Necessity follows by applying $\overline A_i$ to the equality. For sufficiency, define $\overline A_i(z)=\phi(A_i x)$ using any $x$ with $\phi(x)=z$; the condition makes this independent of the choice. Surjectivity proves uniqueness. Iterating the commuting identities gives the corresponding map for every finite word of interventions. If two such words act identically on $X$, their induced maps agree on every $z=\phi(x)$, so this construction respects all relations in the transformation monoid.

\paragraph{Attempt to infer causal adequacy from commutation.}
The condition is too weak on its own. A constant $\phi$ always satisfies it, regardless of the interventions. More informatively, on $X=\{0,1\}^2$ let $A_1$ flip the first bit and $A_2$ flip the second. The projection $\phi(x_1,x_2)=x_1$ induces a flip for $A_1$ and the identity for $A_2$. It preserves commutation but maps distinct interventions $A_2$ and the identity on $X$ to the same intervention on $Z$. It therefore fails faithfulness, in the sense of injectivity on the intervention actions. It also cannot represent the task $o(x)=x_2$.

At the opposite extreme, any bijection $\phi$ has induced maps $\phi A_i\phi^{-1}$ and retains all information. Thus commutation alone does not choose a useful compressed representation between these extremes. Additional task preservation, intervention distinguishability, statistical identifiability, or complexity restrictions must be specified before a canonical choice can be proved.
\section{Verification and adversarial audit}
Surjectivity is used for both definition and uniqueness of induced maps; without it arbitrary behavior off the image would remain. The constant and two-bit examples were checked for both generators and hence for all words. A transformation monoid is sufficient; invertibility or a group is not silently assumed.
\paragraph{Audit corrections.} Explicitly distinguished commuting, faithful, and task-preserving maps; a commuting constant map is not advertised as a satisfactory causal representation.
\paragraph{External computation.} None; symbolic arguments and boundary checks only.
\section{Scope and remaining gaps}
The criterion does not cover stochastic interventions, latent identifiability, do-calculus, or selection of a representation from data. It refines a finite special case of a definition already present in the source, leaving the broad original target unresolved.
\paragraph{Integration review.} Codex (GPT-6; exact serving revision not exposed) reviewed the written arguments and expanded layout and source-reference wrappers. Repository metadata and the production web renderer were checked using Python 3.13 and JavaScript with MathJax 3.2.2. This model-assisted review is not independent human proof verification; no PDF compilation or formal proof checking is asserted.

\section{Completion Estimate}
\noindent\textbf{COMPLETION: 0\%}
The calculations below the preserved source statement investigate a concrete route,
but establish no substantiated new progress toward the entire original target.
Known special cases and failed sufficient conditions are not credited as a solution.
\section{References}
The original source references and catalogue attribution are preserved in the
Sources section above. The following URLs were actually consulted for this attempt;
retention of an additional original reference does not claim fresh consultation of it.
\begin{itemize}
\item \url{https://ems.press/content/serial-article-files/51348} (consulted 2026-10-11). 
\end{itemize}
\section{Model and contributor disclosure}
Codex GPT-6 completed written-proof review of this attempt at its declared scope; this is model-assisted review, not independent human proof checking.
\paragraph{Material model collaboration.} Root Codex GPT-6 supplied the finite deterministic forward-stable fiber criterion and intervention-faithfulness counterexamples. The gpt-6-astra author inspected and reworked these arguments; this is model-assisted mathematical work, not independent human verification.
OpenAI gpt-6-astra, requested reasoning configuration ultra, generated this research
attempt on 2026-10-11. The exact serving revision is not exposed. Codex GPT-6
performed orchestration and repository-format review; its exact serving revision
is also not exposed. The mathematical route, boundary checks and final adversarial
self-audit followed contribution prompts A then B in
\texttt{docs/contribution\_prompts.md}, after reading \texttt{CONTRIBUTING.md}.
Python 3.14.2 assembled the standalone documents; any mathematical computation is
specified in Verification. No theorem prover was used. No human mathematical
authorship, substantial human proof correction, or independent human proof checking
is asserted. Model-assisted review, source retrieval, and structural checks do not
establish correctness or novelty. Full private reasoning transcripts are unavailable
as reproducibility artifacts; the complete arguments, inputs and outcomes actually
used are stated here or in the companion batch audit.
\section{Final status}
Unresolved as a new research contribution. This is a concrete route investigation,
not a claim of new complete resolution or a current-status certification. A known
answer reproduced for context does not become a new contribution. The catalogue's
source-status label is preserved; this attempt status does not assert that a
source-labelled solved question is open.
\end{document}
